% !TEX program = xelatex
\documentclass[12pt,a4paper,fontset=ubuntu]{ctexart}

% ── Packages ──
\usepackage[top=2.5cm, bottom=2.5cm, left=2.5cm, right=2.5cm, headheight=14pt]{geometry}
\usepackage{amsmath}
\usepackage{hyperref}
\usepackage{xcolor}
\usepackage{enumitem}
\usepackage{listings}
\usepackage{tcolorbox}
\usepackage{amssymb}
\usepackage{array}
\usepackage{booktabs}
\usepackage{fancyhdr}
\usepackage{lastpage}
\usepackage{tabularx}

\hypersetup{
  colorlinks=true,
  linkcolor=blue!60!black,
  urlcolor=blue!60!black,
}

% ── Page style ──
\pagestyle{fancy}
\fancyhf{}
\fancyhead[L]{\small RealMan 真机测试操作手册}
\fancyhead[R]{\small 青赋驭境科技}
\fancyfoot[C]{\thepage}

% ── Colors ──
\definecolor{criticalbg}{RGB}{255,235,235}
\definecolor{criticalfg}{RGB}{180,20,20}
\definecolor{checkbg}{RGB}{235,245,235}
\definecolor{infobg}{RGB}{235,240,255}
\definecolor{codebg}{RGB}{245,245,245}
\definecolor{warnbg}{RGB}{255,250,230}

% ── Code listing style ──
\lstset{
  basicstyle=\small\ttfamily,
  backgroundcolor=\color{codebg},
  frame=single,
  framesep=6pt,
  rulecolor=\color{black!20},
  breaklines=true,
  xleftmargin=1em,
  xrightmargin=1em,
  showstringspaces=false,
  aboveskip=8pt,
  belowskip=8pt,
}

% ── Checkbox command ──
\newcommand{\chkbox}{$\square$}
\newcommand{\chkboxdone}{$\boxtimes$}

% ── tcolorbox styles ──
\newtcolorbox{criticalbox}{
  colback=criticalbg,
  colframe=criticalfg,
  sharp corners,
  boxrule=1pt,
  left=4mm, right=4mm, top=3mm, bottom=3mm,
  before skip=10pt, after skip=10pt,
}

\newtcolorbox{infobox}{
  colback=infobg,
  colframe=blue!40,
  sharp corners,
  boxrule=0.5pt,
  left=4mm, right=4mm, top=3mm, bottom=3mm,
  before skip=8pt, after skip=8pt,
}

\newtcolorbox{warnbox}{
  colback=warnbg,
  colframe=orange!60,
  sharp corners,
  boxrule=0.5pt,
  left=4mm, right=4mm, top=3mm, bottom=3mm,
  before skip=8pt, after skip=8pt,
}

\newcommand{\passcriteria}[1]{%
  \textbf{通过标准：} #1%
}

\newenvironment{testitem}[1]{%
  \subsection{#1}
}{%
  \bigskip
}

% ── Title ──
\title{\Huge\textbf{RealMan 真机测试操作手册}\\[6pt]
       \large 密级：内部}
\author{}
\date{\today}

\begin{document}

\maketitle

\vspace{2cm}

\begin{center}
\begin{minipage}{0.7\textwidth}
\large
\textbf{目标：} 按照本手册逐步操作，验证 \texttt{realman\_driver}、
\texttt{realman\_vision}、\texttt{omr\_controller} 三个包
在真机上的功能是否正常。
\end{minipage}
\end{center}

\vfill

\clearpage

\tableofcontents
\clearpage

% ════════════════════════════════════════════════════════════════════
\section{准备工作}

\subsection{硬件检查}

\begin{itemize}[label=\chkbox, itemsep=4pt]
  \item 机械臂已开机，网线已连接到 MiniPC
  \item Intel RealSense D400 相机已通过 USB 3.0 连接到 MiniPC
  \item 棋盘格标定板已打印（$11\times8$ 内角点，方格边长 30\,mm），贴在平整硬板上
  \item 机械臂工作范围内无杂物，急停按钮可触及
\end{itemize}

\subsection{环境初始化}

打开终端，执行以下命令：

\begin{lstlisting}[language=bash]
# 1. 加载 ROS2 环境
source /opt/ros/humble/setup.bash
source ~/ws/realman/install/setup.bash

# 2. 确认机械臂 IP（默认 192.168.1.18，不确定问负责人）
ping 192.168.1.18      # 应该能 ping 通

# 3. 确认相机连接
ls /dev/video*          # 应能看到 realsense 相关 video 设备
\end{lstlisting}

\subsection{重要提醒}

\begin{center}
\begin{tabularx}{\textwidth}{@{}>{\bfseries}lX@{}}
  \toprule
  单位 & 代码中角度用\textbf{弧度}，长度用\textbf{米}。示教器上角度用\textbf{度}，长度用\textbf{毫米}。
        对比时注意换算：$1\,\text{rad} \approx 57.3^\circ$，$1\,\text{m} = 1000\,\text{mm}$。\\[4pt]
  安全 & 运行运动测试前确保周围无人。急停按钮随手可及。\\[4pt]
  报错记录 & 任何步骤报错，截图或复制完整错误信息（包含错误码），记录到测试报告中。\\[4pt]
  测试数据 & 所有测试产物（图片、CSV、YAML）保留，不要删除，测试结束后统一归档。\\
  \bottomrule
\end{tabularx}
\end{center}

% ════════════════════════════════════════════════════════════════════
\section{Phase 1：基础硬件连通性}

\subsection{T1.1 — 机械臂 TCP 连接}

\textbf{操作：}

\begin{lstlisting}[language=bash]
# 运行 gripper_test，自动连接机械臂（默认 IP 192.168.1.18）
ros2 run realman_driver gripper_test
\end{lstlisting}

\textbf{期望看到：}

\begin{lstlisting}[]
[1] Connected to arm
[2] Gripper pre-check:
  enable: 1  online: 1  mode: ...
\end{lstlisting}

\passcriteria{
  不报错，能看到 ``Connected to arm''；夹爪状态 online: 1。
}

\begin{warnbox}
\textbf{如果失败：}
\begin{enumerate}[nosep]
  \item 检查 IP：\texttt{ping 192.168.1.18}
  \item 检查机械臂是否开机（示教器有显示）
  \item 确认机械臂和 MiniPC 在同一网段
  \item 如果 IP 不同，用 \texttt{--ip <实际IP>} 指定
\end{enumerate}
\end{warnbox}

\begin{center}\chkbox\ \textbf{T1.1 通过}\end{center}

% ──

\subsection{T1.2 — RealSense 相机采集}

\textbf{操作：}

\begin{lstlisting}[language=bash]
ros2 run omr_controller collect_data \
    --ip 192.168.1.18 --output /tmp/test_cam --count 5
\end{lstlisting}

\begin{infobox}
有显示器时：弹出相机预览窗口，能看到实时彩色画面。

无显示器时（SSH 远程）：自动降级为 headless 模式，终端输出状态日志。
\end{infobox}

\passcriteria{
  连续能看到画面，不报错；按 \texttt{Q} 可正常退出。
}

\begin{center}\chkbox\ \textbf{T1.2 通过}\end{center}

% ════════════════════════════════════════════════════════════════════
\section{Phase 2：机械臂运动与状态读取}

\subsection{T2.1 — 关节位置读取}

\textbf{操作：}

\begin{lstlisting}[language=bash]
ros2 run realman_driver joint_test
\end{lstlisting}

\begin{infobox}
连接机械臂后每秒打印一次 6 个关节角度（同时显示弧度和度），共 20 次，按 Ctrl+C 可提前退出。
\end{infobox}

\textbf{期望看到：}

\begin{lstlisting}[]
[1] Joints (rad): J1=0.xxx  J2=0.xxx  J3=0.xxx  J4=0.xxx  J5=0.xxx  J6=0.xxx
               (deg): J1=xx.x   J2=xx.x   J3=xx.x   J4=xx.x   J5=xx.x   J6=xx.x
\end{lstlisting}

\passcriteria{
  6 个关节的读数均为正常弧度值（范围 $\approx \pm\pi$）；读数不为全零；
  手动轻推机械臂某关节（松开刹车后），对应读数有变化。
}

\begin{center}\chkbox\ \textbf{T2.1 通过}\end{center}

% ──

\subsection{T2.2 — 末端位姿读取}

\textbf{操作：}

先用示教器记下当前机械臂末端位姿（示教器上通常显示 $x/y/z$ 毫米和
$roll/pitch/yaw$ 度）。

有显示器时，运行采集程序在 overlay 窗口读取：

\begin{lstlisting}[language=bash]
ros2 run omr_controller collect_data --ip 192.168.1.18 --count 1
\end{lstlisting}

\begin{infobox}
有显示器时：CV 窗口左上角显示 ``Arm: x=... y=... z=... roll=... pitch=... yaw=...''
（角度已换算为度）。无显示器时可改用 \texttt{joint\_test}
（输出含 toolPose 但不含此信息，建议通过有显示器的环境测试）。
\end{infobox}

\passcriteria{
  CV overlay 显示的位姿数值与示教器一致。

  \textbf{单位换算：} 示教器角度 $\div 57.3 =$ 弧度，
  示教器位置 $\div 1000 =$ 米。

  \textit{示例：} 示教器显示 $X=500\,\text{mm}$、$Roll=30^\circ$
  $\rightarrow$ overlay 显示接近 $X=0.5\,\text{m}$、$Roll=30^\circ$
  （overlay 已转换为度）。
}

\begin{center}\chkbox\ \textbf{T2.2 通过}\end{center}

% ──

\subsection{T2.3 — \textcolor{red}{moveJ 弧度/度转换验证（最高优先级）}}

\begin{criticalbox}
\textbf{\large [WARNING] 这是最重要的测试。如果转换有 bug，
机械臂会以错误的角度运动，可能造成碰撞。}
\end{criticalbox}

\textbf{背景：} C SDK 内部使用度，API 对外使用弧度。\texttt{Arm::Impl}
在调用 SDK 前做 rad$\rightarrow$deg 转换。\texttt{movej\_test}
正是为此设计——发送 Joint1 $=0.5\,\text{rad}$ 的指令，肉眼判断是否正确。

\textbf{操作：}

\begin{enumerate}[nosep]
  \item 在示教器上记下当前 Joint1 的角度（比如显示 $10^\circ$）
  \item 运行 movej\_test：
\end{enumerate}

\begin{lstlisting}[language=bash]
ros2 run realman_driver movej_test
\end{lstlisting}

\begin{enumerate}[nosep, start=3]
  \item 观察机械臂实际运动后，在示教器上再次读取 Joint1 角度
\end{enumerate}

\passcriteria{
  机械臂 Joint1 实际旋转约 $28.6^\circ$（即 $0.5\,\text{rad}$），
  其他关节基本不动；终端输出 ``MoveJ complete''，不报错。
}

\begin{criticalbox}
\textbf{如果 Joint1 只转了 $\approx0.5^\circ$ 而不是 $28.6^\circ$，
说明弧度$\rightarrow$度转换有 bug！}

立刻停止测试，报告负责人。
\end{criticalbox}

\begin{center}\chkbox\ \textbf{T2.3 通过}\end{center}

% ──

\subsection{T2.4 — 急停功能}

\textbf{操作：}

\begin{lstlisting}[language=bash]
# 终端 1：启动 arm_node
ros2 run realman_driver arm_node --ros-args -p arm_ip:=192.168.1.18

# 终端 2：发送运动指令让机械臂动起来，然后立即急停
ros2 service call /arm_node/stop std_srvs/srv/Trigger "{}"
\end{lstlisting}

\passcriteria{
  机械臂在调用 stop 后 1 秒内完全停止；不报错，服务返回 \texttt{success: true}。
}

\begin{center}\chkbox\ \textbf{T2.4 通过}\end{center}

% ──

\subsection{T2.5 — 直线运动 moveL}

\textbf{操作：}

\begin{lstlisting}[language=bash]
ros2 run realman_driver movel_test
\end{lstlisting}

\begin{infobox}
程序连接机械臂 $\rightarrow$ 读取当前 TCP 位姿
$\rightarrow$ Z 轴 +50\,mm $\rightarrow$ 执行 \texttt{moveL}
$\rightarrow$ 读取并打印最终位姿和 Z 偏移量。
\end{infobox}

\textbf{期望看到：}

\begin{lstlisting}[]
Current TCP: x=... y=... z=... roll=... pitch=... yaw=...
Target TCP:  x=... y=... z=... roll=... pitch=... yaw=...
Executing moveL (speed=20, blocking)...
Final TCP:   x=... y=... z=... roll=... pitch=... yaw=...
Delta Z = 0.0500 m  (expected ~0.050)
moveL test PASSED
\end{lstlisting}

\passcriteria{
  机械臂末端沿直线上升约 50\,mm；示教器上轨迹为直线；终端 Delta Z $\approx 0.050\,\text{m}$。
}

\begin{center}\chkbox\ \textbf{T2.5 通过}\end{center}

% ════════════════════════════════════════════════════════════════════
\section{Phase 3：夹爪操作}

\begin{warnbox}
如果机械臂未安装夹爪，此阶段可跳过。
\end{warnbox}

\subsection{T3.1 — 夹爪基本开合}

\textbf{操作：}

\begin{lstlisting}[language=bash]
ros2 run realman_driver gripper_test
\end{lstlisting}

观察各步骤：步骤 [4] 夹爪完全张开；步骤 [5] 闭合到一半（500）；
步骤 [7] 完全闭合（1000）；步骤 [8] 再次张开。

\passcriteria{
  各步骤夹爪动作正常；终端打印 ``All gripper tests passed.''；无超时报错。
}

\begin{center}\chkbox\ \textbf{T3.1 通过}\end{center}

% ──

\subsection{T3.2 — 力控抓取 gripperPick}

\textbf{操作：}

在夹爪中间放入一个物体（如纸杯、小块泡沫），运行 \texttt{gripper\_test}。

\passcriteria{
  步骤 [6] 夹爪闭合到接触物体后自动停止（不会夹扁物体）；
  gripperState 输出 mode 字段显示力控模式。
}

\begin{center}\chkbox\ \textbf{T3.2 通过}\end{center}

% ──

\subsection{T3.3 — 夹爪状态读取}

\textbf{操作：}

分三次运行 gripper\_test，每次在启动前手动将夹爪调整到不同位置
（全开/半开/全闭）。

\passcriteria{
  三次的 \texttt{actpos} 值递增或递减一致；
  $\text{actpos}=0$ 接近全开，$\text{actpos}=1000$ 接近全闭。
}

\begin{center}\chkbox\ \textbf{T3.3 通过}\end{center}

% ════════════════════════════════════════════════════════════════════
\section{Phase 4：棋盘格检测验证}

\subsection{T4.1 — 棋盘格检出率}

\textbf{操作：}

\begin{lstlisting}[language=bash]
ros2 run omr_controller collect_data \
    --ip 192.168.1.18 --output /tmp/calib_test --count 18
\end{lstlisting}

启动后进入交互模式：
\begin{itemize}[nosep]
  \item \textcolor{green!60!black}{绿色圆点 overlay} = 棋盘格角点被成功检测
  \item \textcolor{red}{红色叉号} = 未检测到
  \item 看到绿色角点 $\rightarrow$ 按 \textbf{S} 保存当前帧
  \item 按 \textbf{Q} 退出
\end{itemize}

将棋盘格放在相机视野的\textbf{至少 10 个不同角度}，每个角度按一次 S：

\begin{center}
\begin{tabular}{@{}ll@{}}
  \toprule
  角度类型 & 示例 \\
  \midrule
  正面 & 棋盘格正对相机 \\
  左侧 $30^\circ$ & 向左倾斜 \\
  右侧 $30^\circ$ & 向右倾斜 \\
  上仰 $30^\circ$ & 向上倾斜 \\
  下俯 $30^\circ$ & 向下倾斜 \\
  左旋 $45^\circ$ & 顺时针旋转 \\
  右旋 $45^\circ$ & 逆时针旋转 \\
  近距 & 棋盘格靠近相机 \\
  远距 & 棋盘格拉远 \\
  \bottomrule
\end{tabular}
\end{center}

\passcriteria{
  10 个不同角度中 $\ge 8$ 个成功检出（显示绿色角点）；角点 overlay 稳定，不闪烁。
}

\begin{warnbox}
\textbf{常见问题：}
\begin{itemize}[nosep]
  \item 棋盘格反光 $\rightarrow$ 调整光源角度
  \item 棋盘格太远 $\rightarrow$ 靠近相机
  \item 一直检测不到 $\rightarrow$ 检查 \texttt{--board-w} / \texttt{--board-h} 是否与打印棋盘格内角点数一致
\end{itemize}
\end{warnbox}

\begin{center}\chkbox\ \textbf{T4.1 通过}\end{center}

% ════════════════════════════════════════════════════════════════════
\section{Phase 5：数据采集集成}

\subsection{T5.1 — 完整采集 18 张图片}

\textbf{操作：}

\begin{lstlisting}[language=bash]
ros2 run omr_controller collect_data \
    --ip 192.168.1.18 \
    --output /tmp/calib_test \
    --count 18 \
    --board-w 11 --board-h 8 --square-size 0.03
\end{lstlisting}

按照 T4.1 的方法，在 18 个不同角度各按一次 S。

\begin{warnbox}
\textbf{注意：} 机械臂每张图之间要有明显不同的旋转（不能只平移棋盘格）。
\end{warnbox}

\passcriteria{
  程序结束时提示 ``Rotation diversity OK''（不是 ``WARNING: rotation diversity insufficient''）；
  \texttt{/tmp/calib\_test/images/} 下有 18 张 jpg 图片；
  \texttt{/tmp/calib\_test/robot\_poses.csv} 有 19 行（1 行 header + 18 行数据）。
}

验证命令：
\begin{lstlisting}[language=bash]
ls /tmp/calib_test/images/ | wc -l       # 应输出 18
wc -l /tmp/calib_test/robot_poses.csv     # 应输出 19
\end{lstlisting}

\begin{center}\chkbox\ \textbf{T5.1 通过}\end{center}

% ──

\subsection{T5.2 — 旋转多样性门控验证}

\textbf{操作：}

重新采集一次，故意只\textbf{在相似角度}拍（机械臂不动，只平动棋盘格）。

\passcriteria{
  程序提示 ``WARNING: rotation diversity insufficient''；
  按 Q 可以正常退出，不会崩溃。
}

\begin{center}\chkbox\ \textbf{T5.2 通过}\end{center}

% ════════════════════════════════════════════════════════════════════
\section{Phase 6：全流水线端到端}

\subsection{T6.1 — run\_pipeline 一键标定}

\textbf{方式 A（推荐）：交互式全流程}

\begin{lstlisting}[language=bash]
ros2 run omr_controller run_pipeline \
    --ip 192.168.1.18 \
    --output /tmp/calib_full \
    --count 18 \
    --mode in_hand \
    --board-w 11 --board-h 8 --square-size 0.03
\end{lstlisting}

\begin{infobox}
程序依次执行：
\textbf{Stage 1} 采集图片
$\rightarrow$ \textbf{Stage 2} 相机内参标定
$\rightarrow$ \textbf{Stage 3} 位姿处理
$\rightarrow$ \textbf{Stage 4} 手眼标定。
Stage 1 是交互采集阶段，与 T5.1 操作相同。
\end{infobox}

\textbf{方式 B（分步执行）：}

\begin{lstlisting}[language=bash]
# 第一步：采集数据
ros2 run omr_controller collect_data \
    --ip 192.168.1.18 --output /tmp/calib_full --count 18

# 第二步：相机内参标定
ros2 run omr_controller calibrate_camera \
    --output /tmp/calib_full --board-w 11 --board-h 8 --square-size 0.03

# 第三步：手眼标定
ros2 run omr_controller compute_hand_eye \
    --output /tmp/calib_full --mode in_hand
\end{lstlisting}

\passcriteria{
  所有阶段不报错；最终产生 \texttt{/tmp/calib\_full/calibration\_result.yaml}。
}

\begin{center}\chkbox\ \textbf{T6.1 通过}\end{center}

% ──

\subsection{T6.2 — 标定结果质量评估}

\textbf{操作：}

\begin{lstlisting}[language=bash]
cat /tmp/calib_full/calibration_result.yaml
\end{lstlisting}

\textbf{逐项检查：}

\begin{center}
\begin{tabularx}{\textwidth}{@{}l>{\ttfamily}lX@{}}
  \toprule
  \textbf{检查项} & \textbf{字段} & \textbf{标准} \\
  \midrule
  旋转矩阵 & rotation\_matrix & $3\times3$，$\det \approx 1.0$（误差 $< 0.01$）\\
  平移向量 & translation\_vector & $3\times1$，值 $< 2\,\text{m}$\\
  重投影误差 & reprojection\_error & $< 100.0$\,px\\
  条件数 & condition\_number & $> 0$ 且 $< 1000$\\
  求解方法 & method & Tsai / Park / Horaud / Daniilidis\\
  目测验证 & 观察 $R$ 矩阵 & 若相机固定在夹爪附近正方向，$R$ 应接近单位阵\\
  \bottomrule
\end{tabularx}
\end{center}

\begin{warnbox}
\textbf{辅助判断：} 如果 $R$ 矩阵和单位阵差很多，或 $t$ 向量与相机到法兰的
物理距离差很多，可能标定有问题。
\end{warnbox}

\begin{center}\chkbox\ \textbf{T6.2 通过}\end{center}

% ════════════════════════════════════════════════════════════════════
\section{Phase 7：ROS2 节点集成}

\subsection{T7.1 — calib\_node 启动}

\textbf{操作：}

\begin{lstlisting}[language=bash]
# 终端 1：启动 calib_node
ros2 run omr_controller calib_node --ros-args -p arm_ip:=192.168.1.18
\end{lstlisting}

验证：

\begin{lstlisting}[language=bash]
# 终端 2：检查节点是否存在
ros2 node list
\end{lstlisting}

\passcriteria{ 输出中能看到 \texttt{/calib\_node}。}

\begin{center}\chkbox\ \textbf{T7.1 通过}\end{center}

% ──

\subsection{T7.2 — 服务调用}

\textbf{操作：}

在 calib\_node 运行的状态下，依次调用各服务：

\begin{lstlisting}[language=bash]
# 相机标定服务（立即返回）
ros2 service call /calib_node/calibrate std_srvs/srv/Trigger "{}"

# 手眼标定服务（立即返回，需先跑过 calibrate）
ros2 service call /calib_node/hand_eye std_srvs/srv/Trigger "{}"

# 紧急停止服务（立即返回）
ros2 service call /calib_node/stop std_srvs/srv/Trigger "{}"
\end{lstlisting}

\begin{warnbox}
\textbf{注意：} \texttt{\~/run} 服务会进入交互采集模式（与 \texttt{collect\_data}
相同），通过 \texttt{ros2 service call} 调用会一直阻塞直到人工按 Q 或 S 完成采集，
不适合作为一键验证。建议用 T6.1 的 \texttt{run\_pipeline} 替代。
\end{warnbox}

\passcriteria{
  每个 service call 返回 \texttt{success: true}；
  calibrate 返回相机标定结果；
  hand\_eye 返回手眼标定结果；
  stop 返回 \texttt{success: true} 且机械臂急停。
}

\begin{center}\chkbox\ \textbf{T7.2 通过}\end{center}

% ──

\subsection{T7.3 — TF 坐标变换广播验证}

\textbf{操作：}

在 calib\_node 完成 hand\_eye 标定后：

\begin{lstlisting}[language=bash]
ros2 run tf2_ros tf2_echo base_link camera_link
\end{lstlisting}

\passcriteria{
  能看到连续的 TF 变换输出（Translation + Rotation）；不报 ``can't find transform''。
}

\begin{center}\chkbox\ \textbf{T7.3 通过}\end{center}

% ════════════════════════════════════════════════════════════════════
\section{测试汇总}

\subsection{完成情况}

\renewcommand{\arraystretch}{1.2}
\begin{center}
\begin{tabularx}{\textwidth}{@{}c>{\raggedright\arraybackslash}p{5.5cm}ccc@{}}
  \toprule
  \textbf{阶段} & \textbf{测试项} & \textbf{状态} & \textbf{备注} \\
  \midrule
  P1 & T1.1 — Arm TCP 连接 & \chkbox & \\
  P1 & T1.2 — 相机采集 & \chkbox & \\\addlinespace
  P2 & T2.1 — 关节位置读取 & \chkbox & \\
  P2 & T2.2 — 末端位姿读取 & \chkbox & \\
  P2 & \textbf{\textcolor{red}{T2.3 — 弧度/度转换}} & \chkbox & \textbf{\textcolor{red}{关键}} \\
  P2 & T2.4 — 急停 & \chkbox & \\
  P2 & T2.5 — 直线运动 moveL & \chkbox & \\\addlinespace
  P3 & T3.1 — 夹爪基本开合 & \chkbox & \\
  P3 & T3.2 — 力控抓取 & \chkbox & \\
  P3 & T3.3 — 夹爪状态读取 & \chkbox & \\\addlinespace
  P4 & T4.1 — 棋盘格检出率 & \chkbox & \\\addlinespace
  P5 & T5.1 — 18 张完整采集 & \chkbox & \\
  P5 & T5.2 — 旋转多样性门控 & \chkbox & \\\addlinespace
  P6 & T6.1 — 全流水线标定 & \chkbox & \\
  P6 & T6.2 — 标定结果质量 & \chkbox & \\\addlinespace
  P7 & T7.1 — calib\_node 启动 & \chkbox & \\
  P7 & T7.2 — 服务调用 & \chkbox & \\
  P7 & T7.3 — TF 广播 & \chkbox & \\
  \bottomrule
\end{tabularx}
\end{center}

\subsection{测试环境记录}

\begin{center}
\begin{tabularx}{\textwidth}{@{}>{\bfseries}lX@{}}
  \toprule
  测试日期 & 202\_\ \rule{1cm}{0.4pt}\ -\ \rule{1cm}{0.4pt}\ -\ \rule{1cm}{0.4pt} \\
  机械臂型号 & \rule{5cm}{0.4pt} \\
  机械臂 IP & \rule{5cm}{0.4pt} \\
  相机型号 & \rule{5cm}{0.4pt} \\
  标定板规格 & 内角点 $\rule{1cm}{0.4pt} \times \rule{1cm}{0.4pt}$，方格 $\rule{1.5cm}{0.4pt}$\,mm \\
  \bottomrule
\end{tabularx}
\end{center}

% ════════════════════════════════════════════════════════════════════
\section*{附录：常见问题}

\addcontentsline{toc}{section}{附录：常见问题}

\subsection*{A. 网络不通}

\begin{lstlisting}[language=bash]
# 确认同一网段
ip addr show       # 看 MiniPC 的 IP
ping 192.168.1.18

# 确认端口 8080 开放
nc -zv 192.168.1.18 8080
\end{lstlisting}

\subsection*{B. 相机无法打开}

\begin{lstlisting}[language=bash]
# 检查 realsense 设备
rs-enumerate-devices

# 重新插拔 USB 线（确认是 USB 3.0 蓝色口）
\end{lstlisting}

\subsection*{C. 棋盘格检测不到}

\begin{enumerate}[nosep]
  \item 确保光线均匀，无强烈反光
  \item 棋盘格占画面 $1/4 \sim 1/2$ 面积
  \item 棋盘格平整，无弯曲
  \item 用 \texttt{--board-w 11 --board-h 8} 确认内角点数正确
\end{enumerate}

\subsection*{D. headless 模式（无显示器）}

SSH 远程登录时，采集程序自动降级为 headless 模式，无法用鼠标键盘交互。
需要模拟按键保存帧。如果完全无交互能力，联系负责人配置远程桌面或使用
\texttt{--mode batch} 模式。

\end{document}
